\subsection{Attori del sistema}
Gli attori onesti inclusi nel sistema utilizzano il protocollo sicuro per completare il processo di votazione, e sono quattro.
\paragraph{User (U):} L'utente è l'entità del sistema che rappresenta l'elettore legittimato a votare. Quindi il suo obiettivo è quello di contattare il Server per inviare il proprio voto, ossia la sua preferenza sul referendum binario, scegliendo tra \textit{"SI"} e \textit{"NO"}.\\
L'utente vuole essere certo che il suo voto sia anonimo, in modo tale che nessun malintenzionato possa compromettere la comunicazione per scoprire la sua preferenza, e si aspetta che il voto resti segreto anche successivamente.\\
Inoltre l'entità in esame confida che il suo voto raggiunga la destinazione e che gli arrivi una conferma dell'avvenuta registrazione del suo voto e che sia integro all'interno del Server. In caso contrario desidera sapere che non è stato registrato a causa di qualsiasi errore del sistema, in modo tale che possa riprovare in seguito.\\
L'utente si aspetta che nessun malintenzionato utilizzi la sua identità per votare al suo posto, altrimenti l'utente perde la possibilità di votare e si compromette il sistema.\\
Quindi tale elettore legittimato deve avere una prova della propria identità, in modo tale da poter dimostrare al Server di essere un utente legittimo.

\paragraph{Server (S):} L'entità Server è incaricata di contare i voti degli utenti per determinare il risultato finale, senza riconoscere la vera identità del votante. L'entità in esame è quindi incaricata di mantenere anonimi i voti, in modo tale che nessun malintenzionato possa scoprire le preferenze degli utenti.\\
Il Server vuole essere certo che i voti che riceve siano autentici, ossia che provengano da utenti legittimi, e che non siano stati alterati durante la trasmissione.\\
L'entità in esame è incaricata di mantenere anonimi i voti, in modo tale che nessun malintenzionato possa scoprire le preferenze degli utenti.\\
L'attore è anche incaricato di registrare correttamente i voti e di garantire che non vengano registrati voti non validi.

\paragraph{Authenticator (A):} L'authenticator è l'entità che si interpone tra lo user che vuole votare e il Server che deve conservare le preferenze degli elettori. Il suo scopo è quello di controllare l'identità dell'entità che desidera votare, quindi controlla l'autenticità dell'utente in modo tale che possa essere considerato un elettore legittimato a partecipare al referendum.\\
Quindi è un vero e proprio proxy di sicurezza che rappresenta l'intermediario nella comunicazione tra lo user e il server per il processo di votazione.

\paragraph{CA :} 


\subsection{Threat model e Asset}


\subsection{Requisiti del sistema e priorità}
\begin{table}[H]
	\centering
	\small
	\renewcommand{\arraystretch}{1.4} % Aumenta l'altezza delle righe per leggibilità
	\begin{tabularx}{\textwidth}{l X c c c}
		\toprule
		\rowcolor{headergray}
		\textbf{Classe} & \textbf{Metodo (Tipo)} & \textbf{Complessità} & $\bm{WMC_{classe}}$ & $\bm{WMC_{classe}}$ \\ 
		\midrule
		
		\textbf{Library} & Library() \textit{\footnotesize (C)} & Low & 1 & 11\\
		& searchBooks(...) \textit{\footnotesize (I)} & Average & 9 \\
		& viewUsers() \textit{\footnotesize (S)} & Low & 1 \\
		\addlinespace[0.5em] \hline \addlinespace[0.5em]
		
		\textbf{Book} & Book(...) \textit{\footnotesize (C)} & Low & 1 & 1 \\
		\addlinespace[0.5em] \hline \addlinespace[0.5em]
		
		\textbf{Reservation} & Reservation(...) \textit{\footnotesize (C)} & Average & 4 & 4 \\
		\addlinespace[0.5em] \hline \addlinespace[0.5em]
		
		\textbf{User} & User(...) \textit{\footnotesize (C)} & Low & 1 & 57 \\
		& register(...) \textit{\footnotesize (M)} & High & 20 \\
		& login(...) \textit{\footnotesize (M)} & Average & 16 \\
		& logout(...) \textit{\footnotesize (M)} & High & 20 \\
		\addlinespace[0.5em] \hline \addlinespace[0.5em]
		
		\textbf{Staff} & Staff() \textit{\footnotesize (C)} & Low & 1 & 102 \\
		& addBook(...) \textit{\footnotesize (M)} & Average & 16 \\
		& updateBook(...) \textit{\footnotesize (M)} & High & 20 \\
		& removeBook(...) \textit{\footnotesize (M)} & Average & 16 \\
		& viewAllReservations() \textit{\footnotesize (S)} & Low & 1 \\
		& approveReservation(...) \textit{\footnotesize (M)} & Average & 16 \\
		& cancelReservation(...) \textit{\footnotesize (M)} & Average & 16 \\
		& sendReminder(...) \textit{\footnotesize (M)} & Average & 16 \\
		\addlinespace[0.5em] \hline \addlinespace[0.5em]
		
		\textbf{Reg. User} & RegisteredUser(...) \textit{\footnotesize (C)} & Low & 1 & 46 \\
		& reserveBook(...) \textit{\footnotesize (M)} & High & 20 \\
		& returnBook(...) \textit{\footnotesize (M)} & Average & 16 \\
		& viewReservedBooks() \textit{\footnotesize (I)} & Average & 9 \\
		\addlinespace[0.5em] \hline \addlinespace[0.5em]
		
		\textbf{Admin} & Administrator() \textit{\footnotesize (C)} & Low & 1 & 53 \\
		& manageUsers(...) \textit{\footnotesize (M)} & High & 20 \\
		& assignRole(...) \textit{\footnotesize (M)} & Average & 16 \\
		& deleteUser(...) \textit{\footnotesize (M)} & Average & 16 \\
		
		\bottomrule
	\end{tabularx}
	\caption{Analisi dettagliata WMC per classi e metodi (Legenda: C=Constructor, S=Selector, M=Modifier, I=Iterator).}
\end{table}

\subsection{Assunzioni del sistema}